\documentclass[11pt]{article}

\usepackage[utf8]{inputenc}
\usepackage[T1]{fontenc}
\usepackage{mathptmx}
\usepackage[margin=1in]{geometry}
\usepackage{graphicx}
\usepackage{amsmath,amssymb}
\usepackage{booktabs}
\usepackage{siunitx}
\usepackage{url}
\usepackage[hidelinks]{hyperref}
\usepackage{titlesec}
\usepackage{caption}
\captionsetup{font=small,labelfont=bf}
\titleformat{\section}{\normalfont\large\bfseries}{\thesection}{0.6em}{}
\titleformat{\subsection}{\normalfont\normalsize\bfseries}{\thesubsection}{0.6em}{}
\setlength{\parskip}{0.35em}

\title{\vspace{-1.5em}\textbf{An Optical Route to Unfrozen Water Content}\\[0.3em]
\large A proposal for the next permafrost spectroscopy paper}
\author{Aviskar Giri, Vasit Sagan, Ajay Kumar Patel,\\
Cagri Gul, Go Iwahana, Tugce Baser}
\date{August 2026}

\begin{document}
\maketitle
\vspace{-1.5em}

\begin{center}
\rule{0.85\textwidth}{0.4pt}
\end{center}

\noindent
\textbf{In one paragraph.} The amount of water that stays liquid in soil below
freezing controls how frozen ground conducts heat, carries load, and releases
carbon. Every accepted way of measuring it requires touching the soil. Liquid
water and ice absorb infrared light at different wavelengths, so a ratio of two
reflectance bands should reveal how the water is divided between the two
phases. We have checked the physics against published optical constants, found
the predicted signal in spectra we already collected at Poker Flat, and
confirmed that no one has published this measurement. We can begin testing it
within weeks using samples already in the laboratory, because those samples
were scanned with the PSR+ but have never been through the HySpex.

\begin{center}
\rule{0.85\textwidth}{0.4pt}
\end{center}

\section{The gap}

A compilation of \num{416} measured soil freezing curves drawn from \num{254}
publications between 1921 and 2021 lists every method the field uses: nuclear
magnetic resonance, time domain reflectometry, calorimetry, dilatometry,
neutron scattering~\cite{Devoie_2022}. All require contact. The closest optical
neighbour is dielectric spectroscopy, which is electrical rather than
optical~\cite{Bittelli_2004}.

To check whether that is a real gap or an impression, we assembled \num{9021}
publications and cross-tabulated sensing method against soil property.
Reflectance spectroscopy is the single largest method in the corpus with
\num{1833} works. Unfrozen water is an active subject with \num{230} works,
growing from seven publications per five-year interval in 2000 to \num{105} in
the interval beginning 2020. Their intersection contains \emph{one} paper
against \num{46.7} expected if the two were independent, and that paper is
about snow. The mid-infrared equivalent is empty.

The same corpus shows the gap is not for want of capability. Optical
spectroscopy already accounts for \SI{53}{\percent} of the soil carbonate
literature, \SI{49}{\percent} of cation exchange capacity, \SI{42}{\percent} of
pH, and \SI{39}{\percent} of organic carbon. Where it is absent is precisely
the physical state variables of frozen ground: ice content, zero of \num{169}
works; unfrozen water, one of \num{179}; thermal conductivity, one of \num{98}.

There is a mechanical reason. Sentinel-2 carries no band between \num{945} and
\SI{1374}{\nano\meter} and none between \num{1374} and \SI{1610}{\nano\meter}.
Landsat carries none between \num{870} and \SI{1610}{\nano\meter}. The
wavelengths where liquid water and ice differ fall in the gaps of the
instruments most of the field uses. The measurement has been invisible to the
standard toolkit.

\section{What the physics says}

Liquid water absorbs near \num{1450} and \SI{1940}{\nano\meter}; ice absorbs
near \num{1500} and \SI{2020}{\nano\meter}~\cite{Warren_2008}. Representing the
illuminated volume as a dry substrate beneath a water film of equivalent
thickness $L$, and mixing the two phases by their liquid fraction $f$, a
two-pass attenuation gives a logarithmic band ratio
\[
I(f) = \ln\frac{R(\lambda_2,f)}{R(\lambda_1,f)},
\qquad
\frac{\partial I}{\partial f} = 2L\left[(\alpha_i-\alpha_w)_{\lambda_2}-(\alpha_i-\alpha_w)_{\lambda_1}\right],
\]
which cancels the substrate and is exactly linear in the quantity we want.
Table~\ref{tab:tiers} gives three usable band pairs. The absorption ordering
between liquid and ice reverses within each pair, which is what makes the ratio
informative.

\begin{table}[h]
\centering\small
\caption{The three band pairs, from published optical constants.}
\label{tab:tiers}
\begin{tabular}{llrrp{6.2cm}}
\toprule
tier & $\lambda$ (nm) & liquid & ice & character \\
     &                & \multicolumn{2}{c}{($\mathrm{mm}^{-1}$)} & \\
\midrule
1 & 1450 / 1500 & 3.149 / 1.883 & 1.831 / 4.550 & strongest signal; \SI{1450}{\nano\meter} sits inside atmospheric water vapour, so laboratory and contact only \\
2 & 1200 / 1250 & 0.126 / 0.111 & 0.070 / 0.130 & about twenty times weaker, but clear of vapour absorption and therefore the form available from orbit \\
3 & 1940 / 2020 & 12.375 / 5.425 & 6.903 / 9.742 & very strong but saturates in wet soil \\
\bottomrule
\end{tabular}
\end{table}

For a realistic silt freezing from all liquid to one quarter liquid, tier one
predicts a \SI{22}{\percent} to \SI{33}{\percent} relative change in
reflectance. The pair at \num{970} and \SI{1030}{\nano\meter} that the visible
and near-infrared literature might suggest gives under \SI{1}{\percent} and is
not usable.

\section{The signal is already in our data}

We applied the index to spectra collected at Poker Flat for the first paper,
which were never gathered with this question in mind.

\begin{itemize}\itemsep2pt
\item Unfrozen active-layer horizons, \num{36} groups: median index
      \num{0.062}. Frozen core groups at or below \SI{-2}{\degreeCelsius},
      seven groups: median \num{-0.264}. Separated at $p = \num{7e-5}$ by
      permutation over sample groups.
\item A saturated peat horizon at \SI{-12.5}{\degreeCelsius} reads on the
      liquid side under tier one, exactly the opacity failure the model
      predicts for very wet organic soil. Tier two recovers it, and separates
      all groups at $p = \num{1e-5}$.
\item Exposed thaw faces measured in the field are indistinguishable from
      unfrozen horizons ($p = \num{0.82}$). Surface melt masks the interior
      state. This is why the measurement must be made under controlled
      conditions rather than in a soil pit.
\end{itemize}

Two further results place the idea at national scale. Across \num{17612} soil
layers of the public laboratory library, the predicted index change exceeds a
three-sigma detection threshold for \SI{97}{\percent} to \SI{99}{\percent} of
layers, and permafrost-order soils give the largest response of any group. The
implied resolution is \SI{0.017}{\cubic\centi\meter\per\cubic\centi\meter} of
volumetric water, comparable to the contact instruments the field currently
relies on.

\section{What we already have}

\begin{table}[h]
\centering\small
\caption{Data in hand. Only the last row requires new fieldwork.}
\label{tab:data}
\begin{tabular}{p{4.3cm}p{4.0cm}p{6.4cm}}
\toprule
holding & extent & what it gives us \\
\midrule
Poker Flat field spectra &
\num{285} active-layer spectra, \num{95} frozen-core spectra with logged
temperature &
the preliminary evidence above \\
\addlinespace
\textbf{Physical samples in the laboratory} &
the 2025 Poker Flat collection, scanned with PSR+, never scanned with HySpex &
\textbf{an experiment available now, without fieldwork} \\
\addlinespace
Core characterization &
\num{406} logged intervals, ten boreholes &
water content, ice content, density, organic content, salinity, cryostructure
for every specimen \\
\addlinespace
Public laboratory library &
\num{76813} mid-infrared and \num{23647} reflective spectra with paired
properties &
national-scale detectability and the population context \\
\addlinespace
Measured freezing curves &
\num{416} curves, \num{869202} points &
independent comparison for anything we retrieve \\
\addlinespace
Ground temperature records &
one station 2003 to 2026, three network sites 2017 to 2026 &
the temperatures the ground actually occupies, which set the experiment \\
\addlinespace
Baseline specimen scans &
\num{12} fabricated specimens, December 2025 &
protocol development and the frost result \\
\addlinespace
August 2026 campaign &
new Poker Flat samples &
independent specimens spanning the natural property range \\
\bottomrule
\end{tabular}
\end{table}

\section{The experiment, and why it can start now}

\subsection{The samples we already hold}

The 2025 Poker Flat samples sit in the laboratory with PSR+ spectra attached
and no HySpex measurements. Running them through the HySpex would produce three
things at once: a cross-sensor comparison between a contact spectroradiometer
and an imaging spectrometer on identical material, imaging rather than point
measurements so that within-sample heterogeneity becomes visible, and a second
independent look at the same specimens whose properties are already
characterized. None of this requires waiting for August.

The one thing we need to establish before designing around them is their
physical state, that is whether they were air dried, kept moist, or kept frozen,
and whether enough material remains per horizon to remold. That determines
whether they serve as intact specimens or as source material at prescribed
density and saturation.

\subsection{The temperature ladder}

The baseline protocol froze specimens on dry ice to \SI{-60}{\degreeCelsius}
and accepted a scan only when the centre thermocouple read at or below
\SI{-20}{\degreeCelsius}. We propose changing this, for a reason that comes out
of the ground records rather than convenience. At \SI{-20}{\degreeCelsius}
almost all the water is already ice, so the liquid fraction sits at its
residual minimum and the index is being read where it carries the least
information. It is also a state the ground rarely occupies: at the two boreal
forest network sites, \SI{86}{\percent} and \SI{88}{\percent} of all sub-zero
time is warmer than \SI{-5}{\degreeCelsius}.

We propose sampling the ladder from the residence-time distribution of the
ground: closely spaced between \SI{0}{\degreeCelsius} and
\SI{-2}{\degreeCelsius} where the soil spends most of its frozen life and the
index changes fastest, sparser out to \SI{-5}{\degreeCelsius}, and a small
number of colder points retained deliberately, because the exposed tundra site
does reach \SI{-20}{\degreeCelsius} and we want that case represented.

\subsection{Frost}

The December baseline compared specimens scanned with surface frost retained
against specimens with the frost removed, and found clear differences across
all three soil types. Our archived field spectra show the same thing from the
other direction: thaw faces read as unfrozen. Surface state dominates the
measurement, and removing frost immediately before acquisition is the single
most important procedural control. The revised protocol already does this.

\subsection{How the two instruments divide the work}

The PSR+ with a contact probe supplies its own illumination and covers
\SIrange{350}{2500}{\nano\meter}, which makes it the reference instrument for
the physics: it sees all three band pairs at high signal-to-noise on a small
spot. The HySpex covers \SIrange{400}{1000}{\nano\meter} at
\SI{1.2}{\nano\meter} and \SIrange{960}{2500}{\nano\meter} at
\SI{1.7}{\nano\meter}, also spanning all three pairs, but as an image. That
lets us map the index across a specimen surface and see whether the phase
partition is spatially uniform, which matters because ice segregation is not
uniform in real soil. Agreement between the two instruments on the same
specimens is itself a result worth reporting, and it is the step that carries
the laboratory measurement toward airborne and orbital imaging spectrometers.

\section{Why this fits IEEE Transactions on Geoscience and Remote Sensing}

The journal has published laboratory-derived physical models of frozen soil
before, in the microwave domain, where a dielectric mixing model parameterizing
unfrozen water became a standard input to sensor forward
models~\cite{Mironov_2004}. Two recent papers in the same family address
exactly our quantity through microwave measurements~\cite{Wu_2022,Wang_2024}.
The optical counterpart does not exist. Framing matters: presented as a soil
study this belongs in a soils journal, but presented as a forward model, an
index, and an instrument requirement it is squarely within the journal's scope.

\section{What we are asking}

\begin{enumerate}\itemsep2pt
\item Agreement that the phase-partition index is the right target for the next
      paper, with the laboratory measurement as the core and the national
      analysis as supporting scale.
\item Access to the 2025 samples and confirmation of their physical state, so
      that HySpex scanning can begin before the August campaign.
\item A decision on the reference measurement for unfrozen water content, since
      that choice sets the apparatus.
\item A view on the one risk we cannot close by analysis. Freezing changes soil
      structure through ice lensing and expansion, which alters scattering
      independently of the phase partition. The band ratio cancels neutral
      changes in brightness but not spectrally structured ones. A short
      measurement at constant total water content would settle it, and we
      propose doing that first.
\end{enumerate}

\bibliographystyle{IEEEtran}
\bibliography{refs}

\end{document}
